\section{Bulgular}

% Her bulgunun basligi IDDIANIN KENDISI. Basliklar icindekiler listesine
% dusuyor, yani okuyucu sadece icindekilere bakarak argumani gorebiliyor.


\begin{bulgu}{Ekonomik Talepler}


    
\end{bulgu}

\begin{bulgu}{Kimlikler}

\end{bulgu}

\begin{bulgu}{Kadın ve Çocuk}
\end{bulgu}


\begin{bulgu}{Kuruluş Anlatısı}
\end{bulgu}


\begin{bulgu}{Erdoğan ve AKP}

[kanıt: rakam, tablo ya da şekil. Yorumu kanıttan sonra bir cümleyle bağla.]
\end{bulgu}

\begin{bulgu}{Antagonist} % erdoğan ve akp hariç antagonistler ama bunu farklı bir biçimde yazmam gerek. antagonisti anlamazlar

2024 yerel seçimleri antagonist kelime oranını neredeyse 3 kat düşürüyor ve bir daha hiç o çizgiye gelmiyor.
diğer düşüşler: 2016 darbe girişimi, 2019 yerel seçimleri, bahçelinin çağırıs
yükselmeler: referandum 2017 ve 2018 genel seçimleri neredeyse 3'er kat artmış ikisi de
\end{bulgu}
